\chapter{Réalisation}

\newpage

\section*{Introduction}
\addcontentsline{toc}{section}{Introduction}


\section{Outils utilisés}

\subsection{Gestion du projet}

\textbf{L'outil Jira}


Durant le processus de développement, nous avons choisi d'adopter la méthodologie du Test Driven Development\footnote{
Le test driven development est un processus de développement logiciel qui repose sur la conversion des exigences logicielles en cas de test avant que le logiciel ne soit entièrement développé. C'est une approche itérative qui combine le développement, l'écriture des tests unitaires et le refactoring. \cite{tdd}
}
(TDD). Pour mettre en œuvre cette approche, nous avons utilisé JUnit 5 et Mockito pour écrire les cas de test nécessaires. \par

70\%


\section*{Conclusion}
\addcontentsline{toc}{section}{Conclusion}

En conclusion de la phase de réalisation, nous avons concrétisé la solution conçue en mettant en place une architecture technique solide et en utilisant les outils et les technologies de développement adaptés. La structure du projet a été soigneusement définie. Les résultats des tests effectués ont confirmé la fiabilité et la performance de la solution. 

